% TODO checklist:
% - [x] Inspected src/mcp/tools/__init__.py for discovery mechanism
% - [x] Inspected src/mcp/runtime.py for @mcp_tool decorator
% - [x] Inspected Q&A.md Q45 for extensibility guide with difficulty ratings
% - [x] Inspected existing tool implementations for patterns
% - [x] Inspected tests/mcp/tools/ for tool testing patterns
% - [x] Cross-referenced with sapthesis-doc.tex subsection on Extensibility Guide

\section{Extending the Tool Ecosystem}
\label{sec:tool-extension}

One of the key design goals of the MCP tools ecosystem is extensibility. The platform provides a well-defined pattern for adding new tools, integrating additional LLM providers, and extending functionality at various layers. This section documents the extension mechanisms and provides practical guidance for developers.

\subsection{The \texttt{@mcp\_tool} Decorator}
\label{subsec:mcp-tool-decorator-usage}

The \texttt{@mcp\_tool} decorator is the primary mechanism for creating new tools. It wraps tool implementations with the runtime scaffolding described in Section~\ref{sec:mcp-runtime}, including authorization, audit logging, telemetry, and error handling.

\subsubsection{Decorator Parameters}

The decorator accepts the following parameters:

\begin{table}[htbp]
    \centering
    \caption{Parameters of the @mcp\_tool decorator}
    \label{tab:mcp-tool-params}
    \begin{tabularx}{\textwidth}{llX}
        \toprule
        \textbf{Parameter} & \textbf{Type} & \textbf{Description} \\
        \midrule
        \texttt{tool\_name} & \texttt{str} & MCP tool name (e.g., ``graph.query''). Must match the module path. \\
        \texttt{required\_scope} & \texttt{str | None} & Required authorization scope. If \texttt{None}, tool is public. \\
        \texttt{rate\_limit} & \texttt{int | None} & Maximum invocations per rate window. If \texttt{None}, no limit. \\
        \texttt{rate\_window} & \texttt{int} & Rate limit window in seconds (default: 60). \\
        \bottomrule
    \end{tabularx}
\end{table}

\subsubsection{Basic Usage}

\begin{lstlisting}[language=Python, caption={Basic @mcp\_tool decorator usage}, label={lst:mcp-tool-basic}]
from src.mcp.runtime import ToolContext, mcp_tool

@mcp_tool(
    tool_name="custom.my_tool",
    required_scope="tools:basic",
    rate_limit=100,
    rate_window=60,
)
async def invoke(
    ctx: ToolContext,
    payload: dict,
    **kwargs
) -> dict[str, Any]:
    """Tool implementation."""
    action = payload.get("action", "default")
    
    if action == "do_something":
        result = await _do_something(payload)
        return {"ok": True, "action": action, "result": result}
    
    raise ValueError(f"Unknown action: {action}")
\end{lstlisting}

\subsection{Adding New Tools}
\label{subsec:adding-new-tools}

Adding a new tool to the ecosystem is straightforward and requires minimal boilerplate.

\subsubsection{Step-by-Step Guide}

\paragraph{Step 1: Create the Tool Module.}
Create a new Python file at the appropriate location under \texttt{src/mcp/tools/}. The module path determines the tool name:

\begin{lstlisting}[language=bash, caption={Tool module location determines tool name}, label={lst:tool-location}]
# For tool name "custom.my_tool":
src/mcp/tools/custom/my_tool.py

# For tool name "graph.new_feature":
src/mcp/tools/graph/new_feature.py
\end{lstlisting}

\paragraph{Step 2: Define the Payload Schema.}
Create a Pydantic model for input validation:

\begin{lstlisting}[language=Python, caption={Tool payload schema definition}, label={lst:tool-schema}]
from pydantic import BaseModel, Field
from typing import Any, Literal

class MyToolPayload(BaseModel):
    """Input schema for custom.my_tool."""
    action: Literal["process", "analyze", "export"] = "process"
    data: dict[str, Any] = Field(default_factory=dict)
    options: dict[str, Any] | None = None
    limit: int = Field(default=100, ge=1, le=10000)
    timeout_ms: int | None = Field(default=None, ge=100, le=300000)
\end{lstlisting}

\paragraph{Step 3: Implement the Tool.}
Implement the tool following the standard patterns:

\begin{lstlisting}[language=Python, caption={Complete tool implementation}, label={lst:tool-implementation}]
"""
MCP Tool: custom.my_tool

Description of what this tool does and when to use it.

Actions:
- process: Process input data
- analyze: Analyze data patterns
- export: Export results in various formats
"""

from __future__ import annotations

from contextlib import suppress
from typing import Any

from pydantic import BaseModel, Field

# Import MCP runtime
from src.mcp.runtime import ToolContext, mcp_tool, validate_payload

# Import logging
with suppress(Exception):
    from src.logging_setup import get_logger
    logger = get_logger(__name__)
if "logger" not in globals():
    import logging
    logger = logging.getLogger(__name__)


# --- Payload Schema ---
class MyToolPayload(BaseModel):
    action: str = "process"
    data: dict[str, Any] = Field(default_factory=dict)
    limit: int = Field(default=100, ge=1, le=10000)


# --- Action Implementations ---
def _act_process(payload: MyToolPayload, ctx: ToolContext) -> dict[str, Any]:
    """Process action implementation."""
    # Check timeout periodically for long operations
    ctx.check_timeout()
    
    result = {"processed": len(payload.data), "items": []}
    # ... processing logic ...
    
    return {
        "ok": True,
        "action": "process",
        "result": result,
        "item_count": len(result["items"]),
    }


def _act_analyze(payload: MyToolPayload, ctx: ToolContext) -> dict[str, Any]:
    """Analyze action implementation."""
    ctx.check_timeout()
    
    analysis = {"patterns": [], "statistics": {}}
    # ... analysis logic ...
    
    return {
        "ok": True,
        "action": "analyze",
        "analysis": analysis,
    }


def _act_export(payload: MyToolPayload, ctx: ToolContext) -> dict[str, Any]:
    """Export action implementation."""
    return {
        "ok": True,
        "action": "export",
        "format": "json",
        "data": payload.data,
    }


# --- Tool Entry Point ---
@mcp_tool(
    tool_name="custom.my_tool",
    required_scope="tools:basic",
    rate_limit=60,
)
def invoke(payload: dict[str, Any] | None = None, **kwargs) -> dict[str, Any]:
    """
    Entry point for custom.my_tool.
    
    Routes to action-specific implementations based on payload.action.
    """
    # Extract context
    ctx = kwargs.get("ctx") or ToolContext(tool="custom.my_tool", action="unknown")
    
    # Validate payload
    validated = validate_payload(payload, MyToolPayload)
    ctx.action = validated.action
    
    # Route to action handler
    action_handlers = {
        "process": _act_process,
        "analyze": _act_analyze,
        "export": _act_export,
    }
    
    handler = action_handlers.get(validated.action)
    if not handler:
        raise ValueError(f"Unknown action: {validated.action}")
    
    return handler(validated, ctx)
\end{lstlisting}

\paragraph{Step 4: Add Unit Tests.}
Create tests in \texttt{tests/mcp/tools/}:

\begin{lstlisting}[language=Python, caption={Tool test implementation}, label={lst:tool-tests}]
# tests/mcp/tools/test_custom_my_tool.py

import pytest
from src.mcp.tools.custom.my_tool import invoke, MyToolPayload
from src.mcp.runtime import ToolContext


@pytest.mark.unit
def test_process_action():
    """Test the process action."""
    result = invoke(
        {"action": "process", "data": {"key": "value"}},
        ctx=ToolContext(tool="custom.my_tool", action="process")
    )
    assert result["ok"] is True
    assert result["action"] == "process"


@pytest.mark.unit
def test_unknown_action_raises():
    """Unknown action should raise ValueError."""
    with pytest.raises(ValueError, match="Unknown action"):
        invoke({"action": "invalid"})


@pytest.mark.unit
def test_payload_validation():
    """Invalid payload should raise ValidationError."""
    with pytest.raises(Exception):  # ValidationError_
        invoke({"limit": -1})  # limit must be >= 1
\end{lstlisting}

\paragraph{Step 5: Verify Auto-Registration.}
The tool is automatically discovered and registered. Verify with:

\begin{lstlisting}[language=Python, caption={Verifying tool auto-registration}, label={lst:verify-registration}]
from src.mcp.tools import list_tools, discover

# Check tool appears in list
tools = list_tools()
assert "custom.my_tool" in tools

# Check discovery returns full spec
specs = discover()
my_spec = next(s for s in specs if s.name == "custom.my_tool")
assert my_spec.callable_name == "invoke"
\end{lstlisting}

\subsection{Extensibility Guide with Difficulty Ratings}
\label{subsec:extensibility-guide}

Table~\ref{tab:extensibility-difficulty} summarizes the difficulty and effort required for common extension tasks.

\begin{table}[htbp]
    \centering
    \caption{Extensibility difficulty ratings for common tasks}
    \label{tab:extensibility-difficulty}
    \begin{tabularx}{\textwidth}{lccX}
        \toprule
        \textbf{Extension Type} & \textbf{Difficulty} & \textbf{Files} & \textbf{Key Steps} \\
        \midrule
        Add new MCP tool & Low & 1--2 & Create module, add tests, auto-registers \\
        Add new LLM provider & Medium & 2--3 & Add adapter, circuit breaker config, health probe \\
        Add new API endpoint & Medium & 2--3 & Create router, schemas, register in app.py \\
        Add new background job & Medium-High & 3--4 & Job handler, worker logic, SSE events \\
        Full feature extension & High & 5+ & Multiple layers, possibly migrations \\
        \bottomrule
    \end{tabularx}
\end{table}

\subsubsection{Adding a New LLM Provider}

Difficulty: \textbf{Medium}

\paragraph{Step 1: Add Provider Type.}
Add the provider to the Provider model:

\begin{lstlisting}[language=Python, caption={Adding provider type to model}, label={lst:add-provider}]
# db/postgres_control/models/provider.py
class ProviderType(str, Enum):
    OPENAI = "openai"
    OLLAMA = "ollama"
    AZURE = "azure"
    ANTHROPIC = "anthropic"  # New provider
\end{lstlisting}

\paragraph{Step 2: Implement Adapter.}
Create adapter in \texttt{src/adapters/}:

\begin{lstlisting}[language=Python, caption={LLM provider adapter implementation}, label={lst:provider-adapter}]
# src/adapters/llm_anthropic.py

import anthropic
from src.config import settings
from src.resilience.circuit import CircuitBreaker

_circuit = CircuitBreaker(name="anthropic", failure_threshold=3)

def complete(
    prompt: str,
    model: str = "claude-3-opus-20240229",
    temperature: float = 0.7,
    max_tokens: int = 1000,
    **kwargs
) -> dict:
    """Anthropic completion via Claude API."""
    with _circuit:
        client = anthropic.Anthropic(api_key=settings.ANTHROPIC_API_KEY)
        response = client.messages.create(
            model=model,
            max_tokens=max_tokens,
            temperature=temperature,
            messages=[{"role": "user", "content": prompt}]
        )
        return {
            "ok": True,
            "content": response.content[0].text,
            "model": model,
            "tokens": response.usage.output_tokens,
        }
\end{lstlisting}

\paragraph{Step 3: Add Health Probe.}
Register health probe for the provider:

\begin{lstlisting}[language=Python, caption={Provider health probe}, label={lst:provider-health}]
# src/health/components.py

async def probe_anthropic() -> ComponentCheck:
    """Check Anthropic API connectivity."""
    try:
        from src.adapters.llm_anthropic import complete
        result = complete("test", max_tokens=5)
        return ComponentCheck(ok=True, status=ComponentStatus.OK)
    except Exception as e:
        return ComponentCheck(ok=False, status=ComponentStatus.ERROR, details={"error": str(e)})
\end{lstlisting}

\subsubsection{Adding a New API Endpoint}

Difficulty: \textbf{Medium}

\paragraph{Step 1: Create Schemas.}

\begin{lstlisting}[language=Python, caption={API endpoint schemas}, label={lst:endpoint-schemas}]
# src/schemas/custom.py

from pydantic import BaseModel, Field

class CustomRequest(BaseModel):
    name: str = Field(..., min_length=1, max_length=100)
    value: int = Field(..., ge=0)

class CustomResponse(BaseModel):
    id: str
    name: str
    value: int
    created_at: str
\end{lstlisting}

\paragraph{Step 2: Create Router.}

\begin{lstlisting}[language=Python, caption={API router implementation}, label={lst:endpoint-router}]
# src/routers/custom.py

from fastapi import APIRouter, Depends, HTTPException
from src.schemas.custom import CustomRequest, CustomResponse
from src.security.deps import get_current_user

router = APIRouter(prefix="/v1/custom", tags=["custom"])

@router.post("/", response_model=CustomResponse, status_code=201)
async def create_custom(
    request: CustomRequest,
    user = Depends(get_current_user),
):
    """Create a new custom resource."""
    # Implementation...
    return CustomResponse(
        id="custom_123",
        name=request.name,
        value=request.value,
        created_at="2025-01-15T10:00:00Z"
    )
\end{lstlisting}

\paragraph{Step 3: Register Router.}

\begin{lstlisting}[language=Python, caption={Registering router in app.py}, label={lst:register-router}]
# src/app.py

from src.routers import custom

def create_app() -> FastAPI:
    app = FastAPI(...)
    # ... other routers ...
    app.include_router(custom.router)
    return app
\end{lstlisting}

\subsubsection{Adding a New Background Job Type}

Difficulty: \textbf{Medium-High}

\paragraph{Step 1: Define Job Handler.}

\begin{lstlisting}[language=Python, caption={Background job handler}, label={lst:job-handler}]
# src/jobs/handlers/custom_job.py

async def handle_custom_job(
    job_id: str,
    payload: dict,
    emit_event: Callable,
) -> dict:
    """Handler for custom job type."""
    total_steps = payload.get("steps", 10)
    
    for i in range(total_steps):
        # Check cancellation
        if await is_cancelled(job_id):
            raise asyncio.CancelledError("Job cancelled")
        
        # Do work
        await process_step(i, payload)
        
        # Emit progress event
        emit_event("progress", {
            "step": i + 1,
            "total": total_steps,
            "percent": ((i + 1) / total_steps) * 100
        })
    
    return {"status": "completed", "steps_processed": total_steps}
\end{lstlisting}

\paragraph{Step 2: Register Handler.}

\begin{lstlisting}[language=Python, caption={Registering job handler}, label={lst:register-handler}]
# src/workers/jobs_worker.py

JOB_HANDLERS = {
    "demo": handle_demo_job,
    "test": handle_test_job,
    "custom": handle_custom_job,  # New handler
}
\end{lstlisting}

\subsection{Testing Tools}
\label{subsec:testing-tools}

Tools should be tested at multiple levels: unit tests, integration tests, and security tests.

\subsubsection{Unit Tests}

Unit tests verify individual action implementations in isolation:

\begin{lstlisting}[language=Python, caption={Unit test structure for tools}, label={lst:unit-test-structure}]
@pytest.mark.unit
class TestMyToolActions:
    """Unit tests for my_tool action implementations."""
    
    def test_process_empty_data(self):
        """Process with empty data returns empty result."""
        result = _act_process(
            MyToolPayload(action="process", data={}),
            ToolContext(tool="custom.my_tool", action="process")
        )
        assert result["ok"]
        assert result["item_count"] == 0
    
    def test_process_with_data(self):
        """Process with data returns processed items."""
        result = _act_process(
            MyToolPayload(action="process", data={"items": [1, 2, 3]}),
            ToolContext(tool="custom.my_tool", action="process")
        )
        assert result["ok"]
        assert result["item_count"] == 3
\end{lstlisting}

\subsubsection{Integration Tests}

Integration tests verify tool behavior with real or simulated dependencies:

\begin{lstlisting}[language=Python, caption={Integration test for tool with database}, label={lst:integration-test}]
@pytest.mark.integration
class TestMyToolIntegration:
    """Integration tests for my_tool with database."""
    
    @pytest.fixture
    def db_session(self):
        """Provide isolated database session."""
        # Setup...
        yield session
        session.rollback()
    
    def test_export_to_database(self, db_session):
        """Export action persists data correctly."""
        result = invoke(
            {"action": "export", "data": {"key": "value"}},
            ctx=ToolContext(tool="custom.my_tool", action="export"),
            db=db_session,
        )
        assert result["ok"]
        # Verify persistence...
\end{lstlisting}

\subsubsection{Security Tests}

Security tests verify authorization and tenant isolation:

\begin{lstlisting}[language=Python, caption={Security tests for tool authorization}, label={lst:security-tests}]
@pytest.mark.security
class TestMyToolSecurity:
    """Security tests for my_tool authorization."""
    
    def test_requires_scope(self):
        """Tool requires tools:basic scope."""
        # Principal without required scope
        ctx = ToolContext(
            tool="custom.my_tool",
            action="process",
            principal={"sub": "user_1", "scopes": []}
        )
        with pytest.raises(PermissionError_):
            invoke({"action": "process"}, ctx=ctx)
    
    def test_tenant_isolation(self):
        """Tool respects tenant boundaries."""
        # Principal from different tenant
        ctx = ToolContext(
            tool="custom.my_tool",
            action="process",
            principal={"sub": "user_1", "scopes": ["tools:basic"], "tenant_id": "tenant_a"},
            tenant="tenant_b"
        )
        with pytest.raises(PermissionError_, match="Tenant mismatch"):
            invoke({"action": "process"}, ctx=ctx)
\end{lstlisting}

\subsection{Best Practices}
\label{subsec:tool-best-practices}

When extending the tool ecosystem, follow these best practices:

\paragraph{Naming Conventions.}
\begin{itemize}
    \item Use lowercase with underscores for module names: \texttt{my\_tool.py}
    \item Use dot notation for tool names: \texttt{category.tool\_name}
    \item Match module path to tool name exactly
\end{itemize}

\paragraph{Action Design.}
\begin{itemize}
    \item Use verbs for action names: \texttt{create}, \texttt{read}, \texttt{process}
    \item Provide a sensible default action
    \item Keep actions focused on single responsibilities
\end{itemize}

\paragraph{Error Handling.}
\begin{itemize}
    \item Raise \texttt{ValidationError\_} for input errors
    \item Raise \texttt{ToolError} for operational errors
    \item Never expose internal details in error messages
    \item Always include error codes for machine parsing
\end{itemize}

\paragraph{Performance.}
\begin{itemize}
    \item Check timeouts periodically in long operations
    \item Implement pagination for large result sets
    \item Use streaming for very large responses
    \item Cache expensive computations appropriately
\end{itemize}

\paragraph{Testing.}
\begin{itemize}
    \item Write unit tests for each action
    \item Include security tests for authorization
    \item Test edge cases and error conditions
    \item Use fixtures for test isolation
\end{itemize}

\subsection{Summary}
\label{subsec:extension-summary}

The MCP tools ecosystem is designed for extensibility:

\begin{itemize}
    \item \textbf{Low-friction tool creation}: Single-file tools with decorator-based registration.
    \item \textbf{Consistent patterns}: All tools follow the same structure and conventions.
    \item \textbf{Automatic discovery}: No manual registration required---tools are discovered by convention.
    \item \textbf{Comprehensive testing}: Unit, integration, and security test patterns are established.
    \item \textbf{Clear difficulty levels}: Extensions range from low (new tool) to high (full feature).
    \item \textbf{Best practices}: Documented conventions ensure maintainability.
\end{itemize}

This extensibility enables the platform to evolve with new requirements while maintaining the governance and observability guarantees of the core framework.
